\section{Discussion}
\label{sec:discussion}
The results support a restrained conclusion. UL-TransXNet is a competitive lightweight ultrasound classifier with strong cross-dataset behavior, but it should not be described as uniformly superior on every dataset and every metric. On TN5000, ConvNeXt-Tiny remains slightly stronger in the repeated-seed table; on AUL, Swin-T slightly exceeds UL-TransXNet in macro F1 and raw accuracy despite lower AUC and balanced accuracy. The proposed model becomes clearly advantageous in the cross-dataset view: it is substantially stronger on BUSI, has higher AUC and balanced accuracy on AUL, and has the best average performance across the three public benchmarks.

The TN5000 ablations clarify the role of the architecture components. The structure-level progression shows that \texttt{TransXNet-GG} is the strongest TN5000 variant, while the \texttt{TransXNet-GGG-MCA} configuration used by UL-TransXNet is the best cross-dataset choice. The module-level ablation shows that differential attention is the clearest contributor to TN5000 thresholded metrics. MUDD has a smaller and less uniform effect. MCA does not uniformly improve TN5000 AUC, but it improves TN5000 F1 and accuracy and produces large BUSI gains. This means that MCA should be framed as a robustness and operating-point module rather than as the sole reason for the model family performance.

This interpretation is important in medical image analysis. Ultrasound datasets differ substantially in organ anatomy, acquisition quality, lesion size, and class prevalence. A small AUC difference on one thyroid benchmark is less informative than the pattern across multiple organs and operating points. The evidence suggests that the proposed architecture is useful because it gives a favorable balance across datasets, not because every module monotonically improves every metric.

The automatic ROI experiment addresses a separate practical concern. The main classifier is trained under an annotation-derived ROI protocol, but the closed-loop BUSI/AUL evaluation shows that detector-predicted crops can replace oracle crops with a moderate performance loss and remain clearly stronger than full-image classification. This result should be interpreted as workflow evidence rather than detector novelty. In particular, AUL localization is weaker than BUSI localization, which explains part of the remaining gap between automatic and oracle ROI classification.

The Android results remain a practical strength of the study. Many ultrasound classification papers stop at workstation-side accuracy. The mobile experiment shows that a TN5000-calibrated prototype can run in a low-latency mode while preserving strong classification quality, and that a more expensive ensemble mode is available when latency is less critical. These measurements do not constitute clinical deployment validation, but they show that UL-TransXNet is compatible with realistic edge-side inference constraints.


\subsection*{Limitations}
This study has several limitations. First, the evaluation is still based on public retrospective datasets, and label definitions, class proportions, and preprocessing conditions are not fully uniform across datasets. Second, the work focuses on binary benign--malignant classification rather than detection, segmentation, or multi-class clinical diagnosis. Third, the automatic ROI experiment uses a practical one-class detector only to form the classifier input; it does not establish a new lesion-detection method, and the AUL results show that localization quality can still limit closed-loop performance. Fourth, the mobile experiment demonstrates feasibility but not a full clinical workflow under realistic operator interaction and acquisition variability. Finally, external robustness remains an open problem because the evaluation is still limited to public retrospective datasets rather than prospective multi-center cohorts. Future studies should therefore focus on broader external validation, stronger detector-classifier integration, calibration under distribution shift, and clinically grounded prospective evaluation.
